\chapter{Design Report Generation}
\label{ch:report}

\emph{Pull requests: \#138, \#145, \#155, \#162, \#169, \#187, \#188.}

\section{Overview}
The single largest part of my internship was the automatically generated \textbf{design
report}---the PDF that summarises an entire plate-girder bridge design, from input parameters
and loads through analysis, design checks, drawings, and quantities. Over seven pull requests I
built and completed multiple chapters, sections, and summary tables of this report, added a
key-first fallback strategy so tables are never left blank, contributed a new design-log
chapter with its supporting logger, made the set of included chapters user-customisable, and
fixed a bug in how report assets were written to disk. This chapter presents that work,
organised the way the report itself is organised.

\section{Report Generator Architecture}
The report is produced by \texttt{core/reports/report\_generator.py}. Each chapter of the
output PDF is built by its own function, which returns a block of LaTeX that the generator
assembles and compiles. The chapter functions are:

\begin{lstlisting}[caption={Per-chapter functions in the report generator (\texttt{report\_generator.py}).},label={lst:reportfns}]
executive_summary()      # overview of the whole design
ch1_project_info()       # project metadata
ch2_input_parameters()   # input configuration
ch3_loads()              # load cases (dead, live, wind, seismic, thermal)
ch4_analysis()           # analysis results
ch5_design_checks()      # design verification tables
ch6_drawings()           # CAD drawings and section previews
ch7_quantities()         # material take-off
ch8_standards()          # standards & assumptions
ch9_design_log()         # NEW: stage-completion log window
ch9_references()         # references
\end{lstlisting}

Because each chapter is an independent function reading from a shared \texttt{output\_dict} and
\texttt{input\_dict}, I could work on one report chapter at a time---which is exactly how the
work was split across my pull requests.

\section{Getting Chapter 6 Working}
My first report contribution (PR~\#138) got Chapter~6 (Drawings) of the generated report
working end-to-end within the generate-report flow, establishing the pattern I would reuse for
the remaining chapters: pull the relevant values and figures from the design output, format them
as LaTeX, and hand the block back to the generator.

\section{Chapter 3 -- Loads and the Key-First Fallback}
Chapter~3 of the report tabulates the load cases---dead, live, wind, seismic, and thermal. The
challenge (PR~\#187) is that some values are supplied directly by the user while others must be
\emph{computed} from IRC~6:2017. Rather than leaving blanks or hard-coding numbers, I adopted a
consistent \emph{key-first fallback}: try to read the value from the input dictionary; only if
it is missing, compute it from the relevant IRC~6 clause.

\begin{lstlisting}[caption={Key-first fallback for wind and seismic values (\texttt{report\_generator.py}, \texttt{ch3\_loads}).},label={lst:fallback}]
# Wind: use the stored hourly-mean values, else compute from IRC6 Table 12
vz_val = input_dict.get(KEY_WL_HOURLY_MEAN_WIND)
pz_val = input_dict.get(KEY_WL_HOURLY_WIND_PRESSURE)
if not vz_val or not pz_val:
    _res = IRC6_2017.table_12(float(_h), _ter, float(_vb))
    if not vz_val:
        vz_val = _res.get("Vz")

# Seismic: use stored coefficients, else compute from IRC6 Cl. 218.5.1
sl_zone_factor = input_dict.get(KEY_SL_ZONE_FACTOR)
sl_ah = input_dict.get(KEY_SL_HORIZONTAL_COEFF)
if not sl_ah or not sl_zone_factor:
    _res = IRC6_2017.cl_218_5_1(zone=f"Zone {_z}", ...)
    sl_zone_factor = sl_zone_factor or _res.get("Z")
    sl_ah = sl_ah or _res.get("Ah")
\end{lstlisting}

The same pattern fills the Chapter~2 input-parameter tables and the Executive Summary, so every
table in these sections is populated either from the user's data or from a code-compliant
computed value---never left empty.

\section{Chapter 5 -- Design Checks}
Chapter~5 is the heart of the report: the design-verification tables. This was the largest
single block of table work, spread over three pull requests.

\subsection{Section 5.1 -- Girder and Stiffener Checks (Tables 5.9--5.16)}
PR~\#155 completed Section~5.1, tables 5.9 through 5.16, covering bearing-stiffener capacity,
deflection, and the controlling design summary. Each row compares a \emph{required} demand
against a \emph{provided} capacity and reports a PASS/FAIL status, with failures highlighted in
red:

\begin{lstlisting}[caption={A design-check row: required vs.\ provided, with pass/fail (\texttt{report\_generator.py}).},label={lst:checkrow}]
_bs_r = bridge.output_dict.get(KEY_SD_BS_R)
def _bs_check(resist_key):
    prov = bridge.output_dict.get(resist_key)
    if prov is None or float(prov) <= 0.0 or _bs_r is None:
        return ("N/A", "N/A", "N/A")
    req_s  = f"{_bs_r} kN"
    prov_s = f"{prov} kN"
    status = "PASS" if float(prov) >= float(_bs_r) else r"\textcolor{red}{FAIL}"
    return (req_s, prov_s, status)
\end{lstlisting}

Deflection checks (Table~5.10) iterate per girder, reading live-load and total deflections and
comparing them against the allowable limits, and the design summary (Table~5.13) selects the
\emph{controlling} check---the one with the highest demand-to-capacity ratio (DCR):

\begin{lstlisting}[caption={Selecting the controlling design check by DCR (\texttt{report\_generator.py}).},label={lst:dcr}]
for _gi, (lbl, _) in enumerate(girder_entries, start=1):
    _live_mm  = bridge.output_dict.get(f"{KEY_SD_DEFL_LIVE}.G{_gi}")
    _total_mm = bridge.output_dict.get(f"{KEY_SD_DEFL_TOTAL}.G{_gi}")
    # ... build the row with values, allowable limit, and status ...

ctrl = max(checks, key=lambda c: c.get("dcr") or 0.0)   # controlling load case
\end{lstlisting}

\subsection{Section 5.2 -- Deck Slab Design (Tables 5.17a--g)}
PR~\#145 added Section~5.2, the deck-slab design, defining and using the keys behind tables
5.17(a) through 5.17(g). These tables lay out the deck-slab design parameters and their
code-based checks in the same required-vs-provided format.

\subsection{Summary Table 5.22}
PR~\#162 completed the Section~5 summary, Table~5.22, which consolidates the individual checks
into a single overview of the design's governing results.

\subsection{Girder side-view fix}
While completing Section~5.1 I also corrected the import path for the girder side-view
stiffener CAD preview used by Chapter~6, pointing it at
\texttt{stiffener\_cad\_preview.StiffenerCadPreviewWidget} so the side-view drawing rendered
correctly.

\section{Chapter 6 -- Drawings (Sections 6.3 and 6.4)}
PR~\#169 completed Chapter~6 by adding Sections 6.3 and 6.4---the cross-bracing and
end-diaphragm layouts. Each is rendered as a labelled figure plus a row of three section
previews (the bracing member and its top and bottom chords), laid out side by side using LaTeX
\texttt{minipage} cells:

\begin{lstlisting}[caption={Cross-bracing and end-diaphragm figures with side-by-side previews (\texttt{report\_generator.py}).},label={lst:secfig}]
cbdia = _sec_fig(fig_paths.get('cb_diagram'), '6.3.1', 'Cross Bracing Layout')
eddia = _sec_fig(fig_paths.get('ed_diagram'), '6.4.1', 'End Diaphragm Layout')

def _sec_cell(...):
    return (r'\begin{minipage}[t]{0.31\textwidth}' + ... + r'\end{minipage}')
\end{lstlisting}

The eight figures involved (a diagram plus three previews for each of the cross-bracing and
end-diaphragm) are captured from the \texttt{TransverseMemberDesign} dialog on the output side
(\texttt{output\_dock.py}) and stored in \texttt{figure\_data}; the plot-capture mechanism this
relies on is described in Chapter~\ref{ch:plots}.

\section{The Asset-Folder Fix}
A bug surfaced when downloading the report as a PDF: an \texttt{assets} folder was being created
inside the user's chosen output directory to hold the report logo, and it was left behind after
export. In PR~\#169 I moved all logo handling inside the \texttt{TemporaryDirectory} used for
LaTeX compilation, so the assets exist only for the duration of the build and are cleaned up
automatically---leaving the user's output directory containing just the finished PDF.

\begin{lstlisting}[caption={Report assets confined to the temporary build directory (\texttt{report\_generator.py}, \texttt{generate\_report}).},label={lst:assetfix}]
with TemporaryDirectory() as tmp_dir:
    tmp_assets = os.path.join(tmp_dir, 'assets')   # logo lives only here
    # ... copy logo into tmp_assets, compile LaTeX, emit final PDF ...
# no assets/ folder left in request.output_dir
\end{lstlisting}

\section{A Customisable Report and the Log-Window Chapter}
The final report pull request (PR~\#188) made the report configurable and added a new chapter.

\subsection{Choosing which chapters to include}
The report-options dialog (\texttt{desktop/ui/dialogs/report\_options.py}) presents the report's
chapters as a checklist. Core chapters---the Executive Summary and Chapters 1--3---are
\emph{locked} and always included; the rest are optional. I added the new
\texttt{"design\_log\_window"} chapter to this list, and the dialog returns the set of selected
sections so the generator emits exactly the chapters the user wants:

\begin{lstlisting}[caption={The report chapter checklist, with locked and optional entries (\texttt{report\_options.py}).},label={lst:reportopts}]
_REPORT_CHAPTERS = [
    # (label, canonical key, locked)
    ("Executive Summary", "executive_summary", True),   # locked
    ("Chapter 1 ...",     "ch1_project_info",  True),   # locked
    # ... chapters 4-10 optional, including:
    ("Design Log Window", "design_log_window", False),
]

def get_checked_leaf_sections(self):
    """Return the canonical keys of the chapters the user selected."""
\end{lstlisting}

\subsection{The design-log chapter}
The new \texttt{ch9\_design\_log} function turns the run's success log---the green
stage-completion messages produced during a design---into a chapter of the report. The log
lines are rendered in monospace, LaTeX special characters are escaped, and a fallback message
is shown when there are no entries:

\begin{lstlisting}[caption={The new design-log chapter (\texttt{report\_generator.py}).},label={lst:designlog}]
def ch9_design_log(log_entries: List[str]) -> str:
    """Render the captured stage-completion (success) log as a chapter."""
    if not log_entries:
        return _chapter_header(...) + "No design log entries"
    # ... escape LaTeX specials, format each line in footnote-size monospace ...
\end{lstlisting}

This relies on a companion enhancement to the logger (\texttt{core/utils/logger.py}), where I
extended \texttt{BridgeLogger} to accumulate success messages during a run and expose them:

\begin{lstlisting}[caption={Capturing success messages in the logger (\texttt{logger.py}).},label={lst:logger}]
class BridgeLogger:
    def analysis_start(self):
        self._success_log = []            # clear at the start of a run

    def get_success_log(self) -> List[str]:
        """Return the green (success) log lines captured during the run."""
        return self._success_log
    # _emit() now appends entries whose level == "success"
\end{lstlisting}

\section{Result}
OsdagBridge produces a complete, IS/IRC-based design report in which every chapter I built is
populated with real, code-compliant data: loads with computed fallbacks, full Section~5 design
checks with pass/fail highlighting, cross-bracing and end-diaphragm drawings, a user-selectable
set of chapters, and a design-log chapter capturing the run. The output directory is left clean
after export.

\osdagfig{figure7.1.png}{A page of the generated design report showing Section~5 design
checks with PASS/FAIL status.}{fig:repchecks}

\osdagfig{figure7.2.png}{The report-options dialog: locked core chapters and
optional chapters, including the new Design Log Window.}{fig:repopts}

\osdagfig{figure7.3.png}{Report Chapter~6 drawings: cross-bracing and end-diaphragm
layouts with section previews.}{fig:repdrawings}

\section{Summary}
Across seven pull requests I built a large share of the OsdagBridge design report: the loads
chapter with a key-first fallback strategy, the complete Section~5 design-check tables, the
Chapter~6 cross-bracing and end-diaphragm drawings, a customisable chapter selection, and a new
design-log chapter with its supporting logger---plus the asset-folder fix that keeps report
exports clean. The plots and section previews these chapters embed are captured through the
mechanism described next, in Chapter~\ref{ch:plots}.
